% !TeX root = ../../Book3.tex
% 纲要标题：### 12.1 平坦性的判据
% 纲要位置：工作记录/计划纲要.md，第 2314 行。
% 纲要细目（写作范围）：
% * 理想张量判据
% * 平坦性的局部性
% * 局部化平坦及一般无挠模反例
% * 复用第二卷第7.3节已经证明的理想张量判据，在本节建立其局部检测形式；随后供第13.4节的完备化应用调用

\section{平坦性的判据}\label{b3:sec:1201}

局部化保持正合性，但一般的标量扩张没有这个保证。平坦性正是要求张量积不把单射变坏。判断时无须逐一检查所有模之间的单射：只检查有限生成理想进入原环的包含就够了。这个一般判据的主要证明见第二卷定理7.3.9；本节经自由模子模归纳给出另一证明，再建立素理想处的局部判据。

\subsection{理想张量判据}\label{b3:subsec:1201:01}

\begin{definition}\label{b3:def:flat-module}
设 \(R\) 为交换环，\(M\) 为 \(R\) 模。如果对任意 \(R\) 模之间的单射 \(U\to V\)，诱导映射 \(U\otimes_RM\to V\otimes_RM\) 仍单射，就称 \(M\) 为\term[pingtan-mo]{平坦模}[flat module]。若交换环同态 \(R\to S\) 使 \(S\) 成为平坦 \(R\) 模，就称它为平坦同态。
\end{definition}

对短正合列 \(0\to U\to V\to V/U\to0\)，张量积总是给出右正合列
\[
U\otimes_RM\longrightarrow V\otimes_RM
\longrightarrow(V/U)\otimes_RM\longrightarrow0.
\]
确实，将 \(V/U\) 的生成元和线性关系代入张量积的定义，得到
\((V/U)\otimes_RM\cong(V\otimes_RM)/\operatorname{im}(U\otimes_RM)\)。
唯一未得到保证的是第一个映射仍为单射；平坦性正是排除这里出现新的核。因此此定义等价于张量函子保持短正合列。自由模平坦，因为与自由模张量就是取若干份直和；平坦模的直和及直和项也平坦，可逐分量检查核。

\begin{theorem}[理想张量判据]\label{b3:thm:ideal-flatness-criterion}
设 \(R\) 为交换环，\(M\) 为 \(R\) 模。则 \(M\) 平坦当且仅当对每个有限生成理想 \(I\subseteq R\)，乘法映射
\[
I\otimes_RM\longrightarrow M,\qquad a\otimes m\longmapsto am
\]
单射。此时该映射把 \(I\otimes_RM\) 识别为 \(IM\)。
\end{theorem}
\begin{proof}
必要性直接应用定义于 \(I\subseteq R\)。为证充分性，先把假设推广到任意理想 \(I\)。即使 \(I\) 不有限生成，一个具体张量按定义仍只能涉及有限多个理想元素。因而核中的张量 \(z=\sum_{i=1}^r a_i\otimes m_i\) 来自有限生成子理想 \(I_0=(a_1,\ldots,a_r)\) 的张量积 \(I_0\otimes_RM\) 中的同一个有限和。该有限和映到 \(M\) 中为零，由假设它在 \(I_0\otimes_RM\) 中已经为零，故 \(z=0\)。

理想是秩一自由模的子模；把秩增加一时，可以用最后一个坐标投影将问题再分出一个理想。由此对 \(n\) 归纳，证明对任意子模 \(N\subseteq R^n\)，映射 \(N\otimes_RM\to M^n\) 单射。\(n=1\) 是刚证的理想情形。投影在 \(N\) 上的像为理想 \(I\)，核为 \(N'\subseteq R^{n-1}\)，所以 \(0\to N'\to N\to I\to0\) 正合。张量后的右正合列
\[
N'\otimes_RM\longrightarrow N\otimes_RM\longrightarrow I\otimes_RM\longrightarrow0
\]
说明：若 \(z\in N\otimes_RM\) 映到 \(M^n\) 为零，则其在 \(I\otimes_RM\) 中的像因理想判据为零，故来自某个 \(z'\in N'\otimes_RM\)。这个 \(z'\) 映到 \(M^{n-1}\) 为零，归纳假设给出 \(z'=0\)，于是 \(z=0\)。无限自由模的情形也成立，因为一个核元素只涉及有限多个 \(N\) 元素，它们的坐标支集包含于同一个有限自由直和项。

最后给定任意包含 \(U\subseteq V\)，它未必带有自由坐标，但可以沿自由模满射拉回到已解决的情形。取满射 \(F\to V\)，其中 \(F\) 自由，记核为 \(K\)，并令 \(G\subseteq F\) 为 \(U\) 的原像。于是 \(K\subseteq G\subseteq F\)，且 \(G/K\cong U\)。若 \(z\in U\otimes_RM\) 映到 \(V\otimes_RM\) 为零，先由右正合性将其提升为 \(w\in G\otimes_RM\)。它在 \(F\otimes_RM\) 中的像来自 \(K\otimes_RM\)，所以减去这个来源在 \(G\otimes_RM\) 中的像后，可使 \(w\) 在 \(F\otimes_RM\) 中的像为零。减去的部分映到 \(U\otimes_RM\) 为零，故仍是原来 \(z\) 的提升。刚证的自由模子模情形给出 \(w=0\)，原来的 \(z\) 因而为零。故任意单射都被保持。
\end{proof}

这个证明只用张量积的右正合性、自由模及有限支集，不需要预先引入 \(\operatorname{Tor}\)。同一判据见 Matsumura 的《Commutative Ring Theory》\cite[\S7, Theorem 7.7, p. 50]{Matsumura1989CommutativeRingTheory}。

给定底环系数 \(a_1,\ldots,a_n\)，先在 \(R^n\) 中考察满足 \(\sum_j a_jb_j=0\) 的系数向量。下面的引理说明，若 \(M\) 平坦，\(M\) 中满足同一关系的元组都可用这些底环系数向量作有限组合得到；关系没有因换到 \(M\) 中而多出无法这样表达的解。

\begin{lemma}[平坦模中的有限关系]\label{b3:lem:flat-relations}
设 \(R\) 为交换环，\(M\) 为平坦 \(R\) 模。若 \(a_j\in R\)、\(m_j\in M\)、\(1\le j\le n\)，满足 \(\sum_{j=1}^n a_jm_j=0\)，则存在有限组 \(y_l\in M\)、\(b_{jl}\in R\)，使
\[
m_j=\sum_l b_{jl}y_l,\qquad \sum_j a_jb_{jl}=0\quad\text{对每个}\;l.
\]
\end{lemma}
\begin{proof}
令 \(\varphi:R^n\to R\) 将第 \(j\) 个基向量送到 \(a_j\)，核为 \(K\)。平坦性使
\(K\otimes_RM\to M^n\to M\) 正合，所以元组 \((m_j)_j\) 是有限和 \(\sum_l v_l\otimes y_l\) 的像。写 \(v_l=(b_{jl})_j\in K\)，便同时得到两组等式。
\end{proof}

\subsection{平坦性的局部性}\label{b3:subsec:1201:02}

\begin{theorem}\label{b3:thm:flatness-local}
设 \(R\) 为交换环，\(M\) 为 \(R\) 模。下列条件等价：\(M\) 在 \(R\) 上平坦；对每个素理想 \(\mathfrak p\subseteq R\)，\(M_{\mathfrak p}\) 在 \(R_{\mathfrak p}\) 上平坦；对每个极大理想 \(\mathfrak m\subseteq R\)，\(M_{\mathfrak m}\) 在 \(R_{\mathfrak m}\) 上平坦。
\end{theorem}
\begin{proof}
若 \(M\) 平坦，将 \(R_{\mathfrak p}\) 模之间的单射 \(U\to V\) 看作 \(R\) 模单射，与 \(M\) 在 \(R\) 上张量后仍单射。因为 \(R\setminus\mathfrak p\) 中的元素已在 \(U\) 上可逆，有自然同构
\[
U\otimes_{R_{\mathfrak p}}M_{\mathfrak p}\cong U\otimes_RM,
\qquad u\otimes(m/s)\longmapsto s^{-1}u\otimes m.
\]
其逆把 \(u\otimes m\) 送到 \(u\otimes(m/1)\)，平衡关系保证两种公式相容。对 \(V\) 也作同样识别，得到局部平坦性。

反向，任取包含 \(U\subseteq V\)，设 \(K\) 为 \(U\otimes_RM\to V\otimes_RM\) 的核。由\cref{b3:thm:localization-exact}及张量积局部化，对每个极大理想有
\[
K_{\mathfrak m}
=\ker\left(U_{\mathfrak m}\otimes_{R_{\mathfrak m}}M_{\mathfrak m}
 \longrightarrow V_{\mathfrak m}\otimes_{R_{\mathfrak m}}M_{\mathfrak m}\right)=0.
\]
\cref{b3:thm:local-zero-detection}于是给出 \(K=0\)，即 \(M\) 平坦。
\end{proof}

\begin{proposition}\label{b3:prop:flat-base-change-transitivity}
设 \(R\to S\) 为交换环同态，\(M\) 为 \(R\) 模，\(N\) 为 \(S\) 模。若 \(M\) 在 \(R\) 上平坦，则 \(S\otimes_RM\) 在 \(S\) 上平坦。若 \(S\) 在 \(R\) 上平坦、\(N\) 在 \(S\) 上平坦，则 \(N\) 在 \(R\) 上平坦。
\end{proposition}
\begin{proof}
第一项对 \(S\) 模包含 \(U\subseteq V\) 使用
\(U\otimes_S(S\otimes_RM)\cong U\otimes_RM\)；右侧的单射由 \(M\) 平坦保证。第二项先与 \(S\) 张量，再与 \(N\) 在 \(S\) 上张量，每一步都保持单射，复合由结合律等于与 \(N\) 在 \(R\) 上张量。
\end{proof}

\subsection{局部化的平坦性与无挠模反例}\label{b3:subsec:1201:03}

局部化 \(S^{-1}R\) 平坦，因为
\(M\otimes_RS^{-1}R\cong S^{-1}M\)，而\cref{b3:thm:localization-exact}已证明右侧函子正合。整环上的平坦模必无挠：对非零 \(a\in R\)，映射 \(R\xrightarrow{\cdot a}R\) 单射，张量后得到乘 \(a\) 在 \(M\) 上单射。

\ProseParagraphBreak
反向通常不成立。设 \(k\) 为域，令 \(R=k[x,y]\)、\(M=I=(x,y)\)，它是整环的子模，因而无挠。乘法使两个生成元满足 \(xy-yx=0\)，因此在 \(I\otimes_RI\) 中将两个乘积分别保留为纯张量，就得到一个可能落入乘法核的元素
\[
z=x\otimes y-y\otimes x\in I\otimes_R I
\]
它在乘法映射下为零；要检验平坦性，关键是这个张量是否已经为零。模掉 \(\mathfrak m=(x,y)\) 后，它的像落在
\((I/\mathfrak mI)\otimes_k(I/\mathfrak mI)\)，其中 \(\bar x,\bar y\) 是基，所列两个不同基张量之差非零，任意特征都如此。因此\cref{b3:thm:ideal-flatness-criterion}说明 \(I\) 不平坦。

\subsection{理想张量判据的局部检测形式}\label{b3:subsec:1201:04}

\begin{corollary}\label{b3:cor:local-ideal-flatness}
设 \(R\) 为交换环，\(M\) 为 \(R\) 模。为证明 \(M\) 平坦，只需对每个有限生成理想 \(I\subseteq R\) 及每个极大理想 \(\mathfrak m\subseteq R\)，证明
\[
I_{\mathfrak m}\otimes_{R_{\mathfrak m}}M_{\mathfrak m}
 \longrightarrow M_{\mathfrak m}
\]
单射。在 Noether 环上，全部理想都有限生成。
\end{corollary}
\begin{proof}
该映射正是 \(I\otimes_RM\to M\) 在 \(\mathfrak m\) 处的局部化。由\cref{b3:thm:local-morphism-detection}，所有这些局部映射单射便使原映射单射，再用\cref{b3:thm:ideal-flatness-criterion}即可。
\end{proof}

这里的检验仍在局部环上进行，不能换成只看剩余域中的维数。设 \(k\) 为域，取 \(R=k[\varepsilon]/(\varepsilon^2)\)、\(\mathfrak m=(\varepsilon)\)、\(M=R/(\varepsilon)\)。环 \(R\) 已是局部环且只有这一个素点，所以茎 \(M_{\mathfrak m}=M\) 仍作为 \(R_{\mathfrak m}=R\) 模保留 \(\varepsilon M=0\) 这条关系。纤维 \(M\otimes_R k\cong M/\mathfrak mM\cong k\) 则是一维 \(k\) 向量空间。

在局部环上，\((\varepsilon)\otimes_RM\cong(\varepsilon)/(\varepsilon)^2\cong k\) 非零，但乘法映射 \((\varepsilon)\otimes_RM\to M\) 将每个张量都送到零，因为 \(\varepsilon M=0\)。因此 \(M\) 不平坦，虽然它的纤维有一维；纤维维数没有记录这条包含映射张量后的非零核。\secref{b3:sec:1203}将用 Fitting 理想记录这个差别。
